% XeLaTeX can use any Mac OS X font. See the setromanfont command below.
% Input to XeLaTeX is full Unicode, so Unicode characters can be typed directly into the source.

% The next lines tell TeXShop to typeset with xelatex, and to open and save the source with Unicode encoding.

%!TEX TS-program = xelatex
%!TEX encoding = UTF-8 Unicode

\documentclass[11pt]{article}


\def\Type{Lesson}
\def\TypeNum{Lesson 12}
\def\TypeTitle{Probability Density Functions}
\def\Semester{Fall 2020}
\def\Course{Math 1220}

\usepackage{preambleDJ} 


% Will Robertson's fontspec.sty can be used to simplify font choices.
% To experiment, open /Applications/Font Book to examine the fonts provided on Mac OS X,
% and change "Hoefler Text" to any of these choices.

%\usepackage{fontspec,xltxtra,xunicode}
%\defaultfontfeatures{Mapping=tex-text}
%\setromanfont[Mapping=tex-text]{Hoefler Text}
%\setsansfont[Scale=MatchLowercase,Mapping=tex-text]{Gill Sans}
%\setmonofont[Scale=MatchLowercase]{Andale Mono}


\begin{document}
\thispagestyle{firstpage}
\topmargin -1cm  %margin looks better starting with second page.
\vspace*{.25in}
{\bf Here's what you said...}\vs{1}\\
{\bf Warmups} 
\enum{
\mpage{.6}{.4}{\vs{-.75}
\item (Breakout rooms) Set up an expression that represents the net volume between $f(x,y)=e^{-x^2-y^2}$ (see  Figure on the right) and the $xy$ plane over the region $R$ that is enclosed by $y=4-(x-2)^2$ and $y=x$ 
}{
\vs{-.75}
\graphicsR{expSurface}{3}
}
\enum{
\mpage{.6}{.4}{
\vs{-1}
\item (Learning Catalytics) Graph the region $R$ on the axes above. A rough graph is fine.\vs{1}
\item Find the intersection points of $y=4-(x-2)^2$ and $y=x$.
}{
\vs{-.5}
\graphicsC{blankAxes}{2}
}
\vs{2}
 \item Set up an expression that represents the net volume over $R$. {\em There's no need to evaluate this expression}.  There are two ways to do this. One approach is easier than the other.
 \newpage
 \item (Challenge question) Set up an expression that represents the {\em average value} of $f$  over the region $R$ Hint: Think about how we defined the average value for  $1$D functions $y=f(x)$ and see if you can extend this idea to $2$D functions $z=f(x,y)$. 
}
\vs{3}
\item (Optional) Given $\ds \int_{-2}^1\int_{y^2}^{2-y} f(x,y)\;dx\;dy$, write down an integral that reverses the order of integration. Hint: start by sketching the region of integration.  \vspace{4in}
}
\newpage
{\bf Probability Density Functions} To understand where we're headed with this, let's look at two examples and then extract the Key Ideas from these examples
\enumR{
\item {\bf Example} Given a typical 6-sided die, what's the probability we roll a $1$?  A $2$? A $6$?
\itemI{
\item {\em Experiment} The act of rolling the die\\
\item {\em Outcomes}  This is the set of all possible results from rolling the die:  $\{ 1,2,3,4,5,6\}$ Note that this set is finitely discrete.\\
\item {\em Random variable} $X$ is the (random) outcome with one roll of the die.  \\
\item {\em Probability function} is denoted by $P$ with domain consisting of the set of outcomes. The range is the probability that the random variable is equal to the input from the domain.  For example 
\enum{
\item What is $P(X= 5)$? \vs{.5}
\item What is $P(x\geq 5)$?\vs{.5}
}
\item {\em Histogram} is a visual representation of each probability when rolling a fair six-sided die.
\graphicsC{1DieHistogram}{2.5}
\vs{1}
{\em Key Ideas} for an experiment with set of outcomes $\{x_1,x_2,x_3\dots,x_n\}$
}
\newpage

%%%%%%%%%%%%%%%%%%sum two dice%%%%%%%%%%%%%%%
\item  (Breakout rooms +Team Learning Catalytics \#1)  The table below contains all possible rolls of two dice and the corresponding sums. You will find it helpful to answer the following questions.
\graphicsC{sum2Dice}{4}
\enum{
\item What is the set of all possible outcomes from adding together the face values of two dice after one roll?  How many different ways can two dice be rolled?   Note that  first rolling a $2$ and then a $1$ is distinct from first rolling a $1$ and then a $2$.   \vspace{1in}
\item What is $P(X=4)$?  It might be helpful to think about how many total combinations of dice rolls will yield a sum of $4$ relative to the total possible number of combinations that lead to our set of outcomes.\vs{.75}\\
\mpage{.6}{.4}{\item What is $P(X\leq 4)$?}{
\item What is $P(\leq 4)+P(X\geq 5)$?}
}
}\vfill
\mpage{.65}{.35}{\hs{1}}{
\graphicsC{sum2DiceHistogram}{2.25}{Sum of 2 dice}
}
\newpage
%%%%%%%%%%%%%%%%%%  Discrete and CTS  PDF  %%%%%%%%%%%%%%%
{\bf Discrete and Continuous Probability Distributions} 
\itemI{
\item A {\em discrete probability distribution} is discrete because the set of outcomes is discrete: In the case of the sum of the roll of 2 dice die we had outcomes $\{2,3,4,\dots,12\}$.  In general, our set of outcomes will be  $\{x_1, x_2, x_3,\dots x_n\}$ where $n$ could also head towards infinity.
\item A {\em continuous probability distribution} is continuous because the set of outcomes of an experiment is continuous.  (e.g. measuring the height of a person, determining the wait time at the DMV, weighing a Roots Bowl, etc).
}
~\\
{\em Some intuition} Let's consider the histograms of the roll of an various numbers of dice.\\
\begin{minipage}{.33\textwidth}
\graphicsC{sum2DiceHistClean}{2}{Sum 2 dice}
\end{minipage}
\hs{.2}
\begin{minipage}{.33\textwidth}
\graphicsC{sum3DiceHistClean}{2}{Sum 3 dice}
\end{minipage}
\hs{.2}
\begin{minipage}{.33\textwidth}
\graphicsC{sum4DiceHistClean}{2}{Sum 4 dice}
\end{minipage}
\begin{minipage}{.33\textwidth}
\graphicsC{sum5DiceHistClean}{2}{Sum 5 dice}
\end{minipage}
\hs{.2}
\begin{minipage}{.33\textwidth}
\graphicsC{expCurveAreaClean}{2}{Sum $\infty$ dice?}
\end{minipage}
\\
%%%%%%%%%%%%%%%%%%  Key Ideas  %%%%%%%%%%%%%%%
{\em Key ideas for continuous random variable $X$ with probability distribution $f(x)$}
\newpage
%%%%%%%%%%%%%%%%%%  Bus stop problem  %%%%%%%%%%%%%%%
\enumR{
\item (Breakout Rooms + Team Learning Catalytics \#2) We're going to determine the probability of catching the bus at the stop on McIntyre between Monroe Hall and the amphitheater.  We assume that a bus comes very $10$ minutes and $T$ is a random variable that represents the wait time to catch the bus.  One probability density function (pdf) for this scenario is shown below. Note that   there's a low probability that a person will have to wait $0$ to $1$ minutes to catch the bus. There's also a low probability that a student will have to wait between $9$ and $10$ minutes to catch the bus.  But, there's a much higher probability that a person will catch the bus in between $4$ and $6$ minutes.  
\graphicsC{tent}{3}
\enum{
\item What's the probability that a person will wait $2.5$ minutes or less? That is, find $P(T\leq 2.5)$.\vs{1}
\item  What's the probability that a person will wait exactly $5$ minutes? That is, find $P(T= 5)$.\vs{1}
\item Let's now assume that it's equally likely that it will take $0-1$ minutes to catch the bus as $1-2$ minutes, as $2-3$ minutes, etc. That is, the wait time is uniformly distributed.  A probability density function that could represent this situation is below. Find a value of $k$ that will guarantee that this is a continuous probability distribution over $0\leq t\leq 10$.
\graphicsL{uniform}{2.5}
}
\newpage
%%%%%%%%%%%%%%%%%%  Extra Problem  %%%%%%%%%%%%%%%
\item (Extra, if time.  Breakout rooms) We're going to make a third set of assumptions to model the wait time for the bus.  Let's assume that there's no guarantee that the buses are running every $10$ minutes. That means it's possible someone could wait, gulp, forever Thus the random variable $T$ is the set of all possible wait times in the interval  $0\leq t<\infty$.  One such pdf is the exponential probability density function $f(t)=2e^{-kt}$ as shown below.\\
\enum{
\mpage{.6}{.4}{
\item Find $k$ that will guarantee that $f$ is a pdf.\vs{1}
}{
\graphicsC{exponential}{3}
}
\vs{2}
\item Find $P(T\leq 1)$\vs{3}
\item What is $P(T>1)$?
}
}
\end{document}

